% GENERATED FILE. Edit canonical lessons, then run npm run generate:books.
% canonical-source-hash: fnv1a64:112f6d222b614193
% canonical-lessons: ML-C360-yadrcchikamayi, ML-C360-manappurvam, ML-C360-abaddham, ML-C360-bhagyam, ML-C360-vidhivasal

\chapter{By Chance or on Purpose}
\label{ch:ml-by-chance-or-on-purpose}
\hlchaptermodality{360}

\begin{chapteropening}
\textbf{By the end of this chapter:} \emph{I can say by chance, on purpose, a blunder, luck and as fate willed.}

Complete the last lesson of chapter 360: I can say by chance, on purpose, a blunder, luck and as fate willed.
\end{chapteropening}

\section[\texorpdfstring{\ml{യാദൃച്ഛികമായി} (yādṛcchikamāyi)}{യാദൃച്ഛികമായി (yādṛcchikamāyi)}]{\ml{യാദൃച്ഛികമായി} (yādṛcchikamāyi) — by chance}

\label{lesson:ML-C360-yadrcchikamayi}



\begin{quote}
Before the new one: say the Malayalam for now and then, then the Malayalam for right after.
\end{quote}

\subsection*{You'll want to know: \ml{യാദൃച്ഛികമായി}}
\textbf{\ml{യാദൃച്ഛികമായി}} — \emph{yādṛcchikamāyi} — \textquotedblleft{}by chance\textquotedblright{}.

\begin{quote}
\ml{ചന്തയിൽ} \ml{വെച്ച്} \ml{ഞാൻ} \ml{യാദൃച്ഛികമായി} \ml{എന്റെ} \ml{പഴയ} \ml{അധ്യാപകനെ} \ml{കണ്ടു}.
\end{quote}

\begin{quote}
\emph{cantayil veccŭ ñān yādṛcchikamāyi enṟe paḻaya adhyāpakane kaṇṭu.}
\end{quote}

\begin{quote}
At the market I met my old teacher by chance.
\end{quote}

The -\ml{ആയി} ending makes an adverb from the adjective \ml{യാദൃച്ഛികം}, chance.

\begin{cousinweb}
\textbf{\ml{യാദൃച്ഛികമായി}} (\emph{yādṛcchikamāyi}) is built from Sanskrit \emph{yādṛcchika}, \textquotedblleft{}spontaneous, accidental, unexpected\textquotedblright{} (Monier-Williams), with the Malayalam adverb ending -\ml{ആയി}.
\end{cousinweb}

\begin{grammarlens}[title={What you've built}]
One more word for telling what happened.
\end{grammarlens}

\subsection*{Your turn}
\begin{itemize}
\raggedright
  \item \emph{Say it:} \emph{yādṛcchikamāyi}
  \item \emph{Say it:} \emph{yādṛcchikamāyi}, once more
  \item \emph{Say it:} \emph{cantayil veccŭ ñān yādṛcchikamāyi enṟe paḻaya adhyāpakane kaṇṭu.} — At the market I met my old teacher by chance.
  \item \emph{Recall:} say the Malayalam for now and then, then the Malayalam for right after, then say \emph{yādṛcchikamāyi} again
\end{itemize}

\begin{tcolorbox}[breakable,colback=teal!4,colframe=teal!35!black,title={Before you move on}]
What does \ml{യാദൃച്ഛികമായി} mean? (\textquotedblleft{}By chance\textquotedblright{}.) Say it once more.
\end{tcolorbox}
\section[\texorpdfstring{\ml{മനപ്പൂർവം} (manappūrvaṁ)}{മനപ്പൂർവം (manappūrvaṁ)}]{\ml{മനപ്പൂർവം} (manappūrvaṁ) — on purpose; deliberately}

\label{lesson:ML-C360-manappurvam}



\begin{quote}
Before the new one: say the Malayalam for right after, then the Malayalam for by chance.
\end{quote}

\subsection*{You'll want to know: \ml{മനപ്പൂർവം}}
\textbf{\ml{മനപ്പൂർവം}} — \emph{manappūrvaṁ} — \textquotedblleft{}on purpose; deliberately\textquotedblright{}.

\begin{quote}
\ml{അവൻ} \ml{മനപ്പൂർവം} \ml{ഗ്ലാസ്} \ml{പൊട്ടിച്ചില്ല}.
\end{quote}

\begin{quote}
\emph{avan manappūrvaṁ glāsŭ poṭṭiccilla.}
\end{quote}

\begin{quote}
He did not break the glass on purpose.
\end{quote}

Literally with the mind first: done with intent. The careful written spelling is \ml{മനഃപൂർവം}.

\begin{cousinweb}
\textbf{\ml{മനപ്പൂർവം}} (\emph{manappūrvaṁ}) is usually traced to Sanskrit \emph{manas}, \textquotedblleft{}mind\textquotedblright{}, and \emph{pūrva}, \textquotedblleft{}preceded by\textquotedblright{}: an act that the mind went before.
\end{cousinweb}

\begin{grammarlens}[title={What you've built}]
One more word for telling what happened.
\end{grammarlens}

\subsection*{Your turn}
\begin{itemize}
\raggedright
  \item \emph{Say it:} \emph{manappūrvaṁ}
  \item \emph{Say it:} \emph{manappūrvaṁ}, once more
  \item \emph{Say it:} \emph{avan manappūrvaṁ glāsŭ poṭṭiccilla.} — He did not break the glass on purpose.
  \item \emph{Recall:} say the Malayalam for right after, then the Malayalam for by chance, then say \emph{manappūrvaṁ} again
\end{itemize}

\begin{tcolorbox}[breakable,colback=teal!4,colframe=teal!35!black,title={Before you move on}]
What does \ml{മനപ്പൂർവം} mean? (\textquotedblleft{}On purpose; deliberately\textquotedblright{}.) Say it once more.
\end{tcolorbox}
\section[\texorpdfstring{\ml{അബദ്ധം} (abaddhaṁ)}{അബദ്ധം (abaddhaṁ)}]{\ml{അബദ്ധം} (abaddhaṁ) — a blunder; a silly mistake}

\label{lesson:ML-C360-abaddham}



\begin{quote}
Before the new one: say the Malayalam for by chance, then the Malayalam for on purpose.
\end{quote}

\subsection*{You'll want to know: \ml{അബദ്ധം}}
\textbf{\ml{അബദ്ധം}} — \emph{abaddhaṁ} — \textquotedblleft{}a blunder; a silly mistake\textquotedblright{}.

\begin{quote}
\ml{ഞാൻ} \ml{ഒരു} \ml{വലിയ} \ml{അബദ്ധം} \ml{ചെയ്തു}.
\end{quote}

\begin{quote}
\emph{ñān oru valiya abaddhaṁ ceytu.}
\end{quote}

\begin{quote}
I made a big blunder.
\end{quote}

By mistake is \ml{അബദ്ധത്തിൽ}: without meaning to.

\begin{cousinweb}
\textbf{\ml{അബദ്ധം}} (\emph{abaddhaṁ}) is borrowed from Sanskrit \emph{abaddha}, \textquotedblleft{}unbound; unmeaning, nonsensical\textquotedblright{} (Monier-Williams). A blunder is something that makes no sense once it is done.
\end{cousinweb}

\begin{grammarlens}[title={What you've built}]
One more word for telling what happened.
\end{grammarlens}

\subsection*{Your turn}
\begin{itemize}
\raggedright
  \item \emph{Say it:} \emph{abaddhaṁ}
  \item \emph{Say it:} \emph{abaddhaṁ}, once more
  \item \emph{Say it:} \emph{ñān oru valiya abaddhaṁ ceytu.} — I made a big blunder.
  \item \emph{Recall:} say the Malayalam for by chance, then the Malayalam for on purpose, then say \emph{abaddhaṁ} again
\end{itemize}

\begin{tcolorbox}[breakable,colback=teal!4,colframe=teal!35!black,title={Before you move on}]
What does \ml{അബദ്ധം} mean? (\textquotedblleft{}A blunder; a silly mistake\textquotedblright{}.) Say it once more.
\end{tcolorbox}
\section[\texorpdfstring{\ml{ഭാഗ്യം} (bhāgyaṁ)}{ഭാഗ്യം (bhāgyaṁ)}]{\ml{ഭാഗ്യം} (bhāgyaṁ) — luck; good fortune}

\label{lesson:ML-C360-bhagyam}



\begin{quote}
Before the new one: say the Malayalam for on purpose, then the Malayalam for a blunder.
\end{quote}

\subsection*{You'll want to know: \ml{ഭാഗ്യം}}
\textbf{\ml{ഭാഗ്യം}} — \emph{bhāgyaṁ} — \textquotedblleft{}luck; good fortune\textquotedblright{}.

\begin{quote}
\ml{ഭാഗ്യത്തിന്} \ml{ഞങ്ങൾക്ക്} \ml{ബസ്} \ml{കിട്ടി}.
\end{quote}

\begin{quote}
\emph{bhāgyattinŭ ñaṅṅaḷkkŭ basŭ kiṭṭi.}
\end{quote}

\begin{quote}
Luckily we caught the bus.
\end{quote}

\ml{ഭാഗ്യത്തിന്} means luckily; \ml{ഭാഗ്യവശാൽ}, fortunately, is built on it.

\begin{cousinweb}
\textbf{\ml{ഭാഗ്യം}} (\emph{bhāgyaṁ}) is borrowed from Sanskrit \emph{bhāgya}, \textquotedblleft{}fortune, one's lot\textquotedblright{} (CDIAL 9434), from \emph{bhaga}, \textquotedblleft{}a share\textquotedblright{}: luck is the share that falls to you.
\end{cousinweb}

\begin{grammarlens}[title={What you've built}]
One more word for telling what happened.
\end{grammarlens}

\subsection*{Your turn}
\begin{itemize}
\raggedright
  \item \emph{Say it:} \emph{bhāgyaṁ}
  \item \emph{Say it:} \emph{bhāgyaṁ}, once more
  \item \emph{Say it:} \emph{bhāgyattinŭ ñaṅṅaḷkkŭ basŭ kiṭṭi.} — Luckily we caught the bus.
  \item \emph{Recall:} say the Malayalam for on purpose, then the Malayalam for a blunder, then say \emph{bhāgyaṁ} again
\end{itemize}

\begin{tcolorbox}[breakable,colback=teal!4,colframe=teal!35!black,title={Before you move on}]
What does \ml{ഭാഗ്യം} mean? (\textquotedblleft{}Luck; good fortune\textquotedblright{}.) Say it once more.
\end{tcolorbox}
\section[\texorpdfstring{\ml{വിധിവശാൽ} (vidhivaśāl)}{വിധിവശാൽ (vidhivaśāl)}]{\ml{വിധിവശാൽ} (vidhivaśāl) — as fate willed}

\label{lesson:ML-C360-vidhivasal}



\begin{quote}
Before the new one: say the Malayalam for a blunder, then the Malayalam for luck.
\end{quote}

\subsection*{You'll want to know: \ml{വിധിവശാൽ}}
\textbf{\ml{വിധിവശാൽ}} — \emph{vidhivaśāl} — \textquotedblleft{}as fate willed\textquotedblright{}.

\begin{quote}
\ml{വിധിവശാൽ} \ml{അന്ന്} \ml{അവൻ} \ml{ആ} \ml{ബസിൽ} \ml{കയറിയില്ല}.
\end{quote}

\begin{quote}
\emph{vidhivaśāl annŭ avan ā basil kayaṟiyilla.}
\end{quote}

\begin{quote}
As fate would have it, he did not get on that bus that day.
\end{quote}

Built like \ml{ഭാഗ്യവശാൽ}, fortunately: by the power of fate.

\begin{cousinweb}
\textbf{\ml{വിധിവശാൽ}} (\emph{vidhivaśāl}) is borrowed whole from Sanskrit \emph{vidhivaśāt}, \textquotedblleft{}through the power of fate or destiny\textquotedblright{} (Monier-Williams), from \emph{vidhi}, \textquotedblleft{}rule, fate\textquotedblright{} (CDIAL 11755), and \emph{vaśa}, \textquotedblleft{}power\textquotedblright{}.
\end{cousinweb}

\begin{grammarlens}[title={What you've built}]
One more word for telling what happened.
\end{grammarlens}

\subsection*{Your turn}
\begin{itemize}
\raggedright
  \item \emph{Say it:} \emph{vidhivaśāl}
  \item \emph{Say it:} \emph{vidhivaśāl}, once more
  \item \emph{Say it:} \emph{vidhivaśāl annŭ avan ā basil kayaṟiyilla.} — As fate would have it, he did not get on that bus that day.
  \item \emph{Recall:} say the Malayalam for a blunder, then the Malayalam for luck, then say \emph{vidhivaśāl} again
\end{itemize}

\begin{tcolorbox}[breakable,colback=teal!4,colframe=teal!35!black,title={Before you move on}]
What does \ml{വിധിവശാൽ} mean? (\textquotedblleft{}As fate willed\textquotedblright{}.) Say it once more.
\end{tcolorbox}
